% Part: incompleteness
% Chapter: incompleteness-provability
% Section: tarski-thm

\documentclass[../../../include/open-logic-section]{subfiles}

\begin{document}

\olfileid[id]{inc}{inp}{tar}

\olsection{Ketakterdefinisian Kebenaran}

Gagasan \emph{keterdefinisian} bergantung pada adanya semantik formal bagi
bahasa aritmetika. Kita telah menguraikan suatu himpunan formula dan kalimat
dalam bahasa aritmetika. ``Interpretasi yang dimaksudkan'' ialah membaca
kalimat-kalimat tersebut sebagai pernyataan mengenai bilangan asli; suatu
pernyataan semacam itu dapat benar atau salah. Misalkan $\Struct{N}$ adalah
!!{structure} dengan domain~$\Nat$ dan interpretasi standar bagi simbol-simbol
dalam bahasa aritmetika. Maka $\Sat{N}{!A}$ berarti ``$!A$ benar dalam
interpretasi standar.''

\begin{defn}
Suatu relasi $R(x_1,\dots,x_k)$ bilangan asli \emph{dapat didefinisikan}
dalam $\Struct{N}$ jika dan hanya jika terdapat formula
$!A(x_1,\dots,x_k)$ dalam bahasa aritmetika sedemikian sehingga untuk setiap
$n_1,\dots,n_k$, $R(n_1,\dots,n_k)$ jika dan hanya jika
$\Sat{N}{!A(\num n_1,\dots,\num n_k)}$.
\end{defn}

Dengan kata lain, suatu relasi dapat didefinisikan dalam $\Struct{N}$ jika
dan hanya jika relasi itu dapat direpresentasikan dalam teori $\Th{TA}$,
dengan $\Th{TA} = \Setabs{!A}{\Sat{N}{!A}}$ sebagai himpunan kalimat
aritmetika yang benar. (Jika hal ini tidak langsung jelas bagi Anda,
kembalilah memeriksa definisi-definisi tersebut dan yakinkan diri Anda bahwa
memang demikian.)

\begin{lem}
Setiap relasi yang dapat dihitung dapat didefinisikan dalam~$\Struct{N}$.
\end{lem}

\begin{proof}
Mudah diperiksa bahwa formula yang merepresentasikan suatu relasi dalam
$\Th{Q}$ mendefinisikan relasi yang sama dalam $\Struct{N}$.
\end{proof}

Sekarang kita dapat bertanya: apakah konversnya juga benar? Dengan kata lain,
apakah setiap relasi yang dapat didefinisikan dalam~$\Struct{N}$ dapat
dihitung? Jawabannya tidak. Misalnya:

\begin{lem}
Relasi penghentian dapat didefinisikan dalam $\Struct{N}$.
\end{lem}

\begin{proof}
Ingat bahwa teorema bentuk normal Kleene menyatakan bahwa setiap fungsi
parsial yang dapat dihitung~$f$ mempunyai indeks~$e$ sedemikian sehingga
$f(x) \simeq U(\umin{s}{T(e,x,s)})$ untuk semua $x \in \Nat$, dengan $U$
dan $T$ rekursif primitif dan karenanya total. Jadi, $f(x)$ terdefinisi
(yakni komputasinya berhenti) jika dan hanya jika terdapat~$s$ sedemikian
sehingga $T(e,x,s)$ berlaku.

Sekarang misalkan $H$ adalah relasi penghentian, yakni
\[
H = \Setabs{\tuple{e,x}}{\lexists[s][T(e, x, s)]}.
\]
Misalkan $!D_T$ mendefinisikan $T$ dalam $\Struct{N}$. Maka
\[
H = \Setabs{\tuple{e,x}}{\Sat{N}{\lexists[s][!D_T(\num e, \num x, s)]}},
\]
sehingga $\lexists[s][!D_T(z, x, s)]$ mendefinisikan~$H$ dalam
$\Struct{N}$.
\end{proof}

\begin{prob}
Tunjukkan bahwa $Q(n) \defiff n \in
\Setabs{\Gn{!A}}{\Th{Q} \Proves !A}$ dapat didefinisikan dalam aritmetika.
\end{prob}

Bagaimana dengan $\Th{TA}$ sendiri? Apakah teori itu dapat didefinisikan
dalam aritmetika? Dengan kata lain, apakah himpunan
$\Setabs{\Gn{!A}}{\Sat{N}{!A}}$ dapat didefinisikan dalam aritmetika?
Teorema Tarski menjawabnya secara negatif.

\begin{thm}
\ollabel{thm:tarski}
Himpunan !!{sentence}s aritmetika yang benar tidak dapat didefinisikan dalam
aritmetika.
\end{thm}

\begin{proof}
Andaikan $!D(x)$ mendefinisikannya, yakni $\Sat{N}{!A}$ jika dan hanya jika
$\Sat{N}{!D(\gn{!A})}$. Berdasarkan lema titik tetap, terdapat suatu formula
$!A$ sedemikian sehingga
$\Th{Q} \Proves !A \liff \lnot !D(\gn{!A})$, dan karena itu
$\Sat{N}{!A \liff \lnot !D(\gn{!A})}$. Namun, kemudian $\Sat{N}{!A}$ jika
dan hanya jika $\Sat{N}{\lnot !D(\gn{!A})}$, yang bertentangan dengan fakta
bahwa $!D(y)$ seharusnya mendefinisikan himpunan pernyataan aritmetika yang
benar.
\end{proof}

Tarski menerapkan analisis ini pada gagasan filosofis yang lebih umum tentang
kebenaran. Untuk sebarang bahasa $L$, Tarski berpendapat bahwa suatu gagasan
kebenaran yang memadai bagi $L$ harus memenuhi, untuk setiap kalimat $X$,
\begin{quote}
`$X$' benar jika dan hanya jika $X$.
\end{quote}
Contoh Tarski yang sering dikutip, untuk bahasa Inggris, adalah kalimat
\begin{quote}
`Salju berwarna putih' benar jika dan hanya jika salju berwarna putih.
\end{quote}
Namun, bagi sebarang bahasa yang cukup kuat untuk merepresentasikan fungsi
diagonal dan sebarang predikat linguistik $T(x)$, kita dapat membangun suatu
kalimat $X$ yang memenuhi ``$X$ jika dan hanya jika bukan
$T(\text{`$X$'})$.'' Karena kita tidak menghendaki predikat kebenaran
menyatakan beberapa kalimat benar sekaligus salah, Tarski menyimpulkan bahwa
kita tidak dapat menentukan predikat kebenaran bagi semua kalimat dalam
suatu bahasa tanpa, dengan suatu cara, melangkah ke luar batas bahasa
tersebut. Dengan kata lain, predikat kebenaran bagi suatu bahasa tidak dapat
didefinisikan dalam bahasa itu sendiri.

\end{document}
